\documentclass[UTF8,11pt]{ctexart}
\usepackage[a4paper,margin=24mm]{geometry}
\usepackage{amsmath,amssymb,amsthm,mathtools}
\usepackage[all]{xy}
\usepackage{xurl,hyperref,enumitem}
\hypersetup{hidelinks,breaklinks=true}
\setlength{\emergencystretch}{3em}
\newtheorem*{theorem}{Theorem}
\newtheorem*{lemma}{Lemma}
\newtheorem*{proposition}{Proposition}
\newtheorem*{corollary}{Corollary}
\theoremstyle{definition}
\newtheorem*{definition}{Definition}
\newtheorem*{remark}{Remark}
\DeclareMathOperator{\Hom}{Hom}
\DeclareMathOperator{\Ext}{Ext}
\DeclareMathOperator{\Tor}{Tor}
\DeclareMathOperator{\Spec}{Spec}
\DeclareMathOperator{\supp}{supp}
\DeclareMathOperator{\im}{im}
\DeclareMathOperator{\id}{id}
\DeclareMathOperator{\Tr}{Tr}
\DeclareMathOperator{\rank}{rank}
\newcommand{\R}{\mathbb R}
\newcommand{\C}{\mathbb C}
\newcommand{\Z}{\mathbb Z}
\newcommand{\N}{\mathbb N}
\newcommand{\Q}{\mathbb Q}
\begin{document}
\section*{015\quad 平方为单位元的群元素计数与Frobenius--Schur指标}
\subsection*{来源与收录范围}
Pavel Etingof, Oleg Golberg, Sebastian Hensel, Tiankai Liu, Alex Schwendner, Dmitry Vaintrob, Elena Yudovina, \emph{Introduction to Representation Theory}；MIT OpenCourseWare 18.712, Fall 2010。使用课程所附 Chapter 4, \emph{Representations of finite groups: further results} 的 PDF，不使用另行出版的 AMS 图书。
\par\url{https://ocw.mit.edu/courses/18-712-introduction-to-representation-theory-fall-2010/84358595a02a73bced2c4e363a5d66f0_MIT18_712F10_ch4.pdf}
\par\url{https://ocw.mit.edu/courses/18-712-introduction-to-representation-theory-fall-2010/pages/lecture-notes/}
\par 获取日期：2026-09-28。\par 位置：\S4.1 的类型定义、Definition 4.3、Theorem 4.4 和完整证明，分章 PDF 第1--2页。
\subsection*{作者命题与人类证明}

Suppose that $G$ is a finite group and $V$ is an irreducible representation of $G$ over $\C$. We say that $V$ is
\begin{itemize}
\item of complex type, if $V\not\cong V^*$,
\item of real type, if $V$ has a nondegenerate symmetric form invariant under $G$,
\item of quaternionic type, if $V$ has a nondegenerate skew form invariant under $G$.
\end{itemize}
\begin{definition}[4.3]
The Frobenius--Schur indicator $FS(V)$ of an irreducible representation $V$ is $0$ if it is of complex type, $1$ if it is of real type, and $-1$ if it is of quaternionic type.
\end{definition}
\begin{theorem}[4.4, Frobenius--Schur]
The number of involutions (=elements of order $\leq2$) in $G$ is equal to $\sum_V\dim(V)FS(V)$, i.e., the sum of dimensions of all representations of $G$ of real type minus the sum of dimensions of its representations of quaternionic type.
\end{theorem}
\begin{proof}
Let $A:V\to V$ have eigenvalues $\lambda_1,\lambda_2,\ldots,\lambda_n$. We have
\begin{align*}
\Tr|_{S^2V}(A\otimes A)&=\sum_{i\leq j}\lambda_i\lambda_j,\\
\Tr|_{\wedge^2V}(A\otimes A)&=\sum_{i<j}\lambda_i\lambda_j.
\end{align*}
Thus,
\[
\Tr|_{S^2V}(A\otimes A)-\Tr|_{\wedge^2V}(A\otimes A)
=\sum_{1\leq i\leq n}\lambda_i^2=\Tr(A^2).
\]
Thus for $g\in G$ we have
\[\chi_V(g^2)=\chi_{S^2V}(g)-\chi_{\wedge^2V}(g).\]
Therefore,
\begin{align*}
|G|^{-1}\chi_V\left(\sum_{g\in G}g^2\right)
&=\chi_{S^2V}(P)-\chi_{\wedge^2V}(P)\\
&=\dim(S^2V)^G-\dim(\wedge^2V)^G\\
&=\begin{cases}
1,&\text{if $V$ is of real type},\\
-1,&\text{if $V$ is of quaternionic type},\\
0,&\text{if $V$ is of complex type}.
\end{cases}
\end{align*}
Finally, the number of involutions in $G$ equals
\[
\frac1{|G|}\sum_V\dim V\,\chi_V\left(\sum_{g\in G}g^2\right)
=\sum_{\text{real }V}\dim V-\sum_{\text{quat }V}\dim V.
\]
\end{proof}

\subsection*{整理说明（非作者正文）}
本题的 involutions 依作者定义包括单位元（阶不超过2），不是只数阶恰好为2的元素。求和遍历不可约复表示同构类。原文 $P$ 是平均投影 $|G|^{-1}\sum_{g\in G}g$；特征标在线性组合上按线性延拓解释。这些是整理者的记号说明。

允许的标准先修为 Schur 引理、平均投影与不变子空间、正规表示特征标和对称/外平方。要求独立建立平方映射的特征标与两个平方表示迹的关系，并返回群元素计数；并非直接应用已知指标求和公式。此题属于资格考试及以上的结构性证明样本，不以篇幅或课程层次标签推定难度。

Problem 4.1 的实代数分类及实结构等价问题不是本题结论，未当作另题或伪称已取到证明。只保留实际使用的类型定义。两页均实际查看原图；没有下载 PDF 原件，包内来源文件为转录。
\subsection*{可修改的简短标注（非作者正文）}
$c$：有限群中平方为单位元的元素计数

$a$：特征标；对称平方；外平方；平均投影；不变双线性型；正规表示

\texttt{cross\_theory\_status = yes}。将平方映射的迹写成对称平方与外平方特征标之差，用不变型区分三种指标，再以正规表示特征标将代数迹返回群元素计数。位置：Theorem 4.4 的全证。

全部标注均为非穷尽、可修改初稿；未标注不表示未使用。
\subsection*{许可与修改记录}
原文来自 MIT OpenCourseWare，作者与课程如上。按来源网站标示的 Creative Commons Attribution--NonCommercial--ShareAlike 4.0 许可复用；本转录及整理沿用同一许可。2026-09-28：助手转录、加入来源定位与简短标注；不暗示作者或 MIT 认可本次整理。许可：\url{https://creativecommons.org/licenses/by-nc-sa/4.0/}。
\end{document}
